\chapter{Conclusões}

\section{Considerações Finais}
\markright{Conclusões}
\label{sec:consideracoes_finais}

O objetivo deste trabalho foi desenvolver e validar tecnicamente uma interface musical digital multissensorial capaz de transformar informações do gesto no controlador Akai MPD218 em comandos para um motor vibrotátil e um LED incorporados a uma unidade vestível. Os resultados permitem responder afirmativamente à questão de viabilidade funcional proposta na Introdução, nas condições de implementação e ensaio descritas.

A cadeia composta por aquisição MIDI, processamento em Python, transmissão UDP sobre Wi-Fi, controle pelo ESP32 e acionamento das saídas foi verificada em etapas e em operação integrada. As simulações fundamentaram o circuito de proteção e o uso do PWM; os ensaios controlados caracterizaram relações mensuráveis entre altura de queda e \textit{velocity}, bem como entre PWM e frequência acústica associada ao motor. A alimentação por baterias permitiu a operação da unidade vestível sem conexão USB com o computador.

A caracterização experimental também indicou a necessidade de rever o mapeamento inicial que emprega referência superior de \textit{velocity} igual a 30, pois essa escolha antecipa a saturação dos comandos. O mapeamento de faixa completa utilizado nos ensaios constitui uma alternativa documentada, ainda não submetida a uma avaliação perceptiva.

A principal contribuição reside na arquitetura implementada e no conjunto de procedimentos técnicos disponibilizados para sua reprodução. O sistema representa ocorrência, duração e intensidade relativa de eventos MIDI, sem converter o conteúdo sonoro em uma reprodução tátil ou luminosa equivalente. A viabilidade demonstrada não permite concluir que a interface melhore a percepção musical de pessoas surdas ou com deficiência auditiva; essa avaliação exige participação da população-alvo em estudos específicos.

\section{Limitações do Protótipo}
\markright{Conclusões}
\label{sec:limitacoes_prototipo}

A principal limitação da caracterização vibrotátil é a ausência de medição inercial direta. O registro com microfone permitiu acompanhar a frequência acústica associada ao motor, mas não forneceu valores de aceleração, deslocamento ou amplitude de vibração. A relação entre PWM e intensidade mecânica transmitida à pele não foi quantificada.

A sequência do ensaio acústico foi executada sempre em ordem crescente de PWM. Variações entre repetições podem incluir efeitos de tempo de funcionamento, aquecimento, condição da alimentação ou acoplamento mecânico e acústico. A correlação observada sustenta a tendência entre comando e frequência registrada, mas não separa essas fontes de variabilidade.

O canal visual foi validado funcionalmente, sem medição fotométrica ou estudo de equivalência perceptiva com o canal tátil. O uso do mesmo valor PWM representa apenas uma estratégia comum de controle.

A comunicação UDP não foi caracterizada quanto à latência fim a fim, e os testes não abrangeram diferentes redes, distâncias ou condições de interferência. A alimentação embarcada também não foi avaliada em termos de autonomia, consumo médio, eficiência do conversor ou temperatura em operação prolongada.

A estrutura vestível foi utilizada como suporte experimental. Conforto prolongado, distribuição de massa, pressão de contato, estabilidade durante movimentos e posicionamento ótimo do motor não foram avaliados por protocolo específico.

A configuração física utilizou um único motor vibrotátil. Não foram avaliados arranjos com múltiplos atuadores, codificação espacial de eventos ou distribuição dos estímulos entre regiões do corpo.

Não foram realizados ensaios formais com pessoas surdas ou com deficiência auditiva. O trabalho não determina limiares de percepção, capacidade de discriminação entre níveis de estímulo, compreensão de padrões rítmicos, conforto de uso ou contribuição efetiva para a experiência musical. A validação é estritamente técnica.

\section{Propostas de Continuidade}
\markright{Conclusões}
\label{sec:propostas_continuidade}

A continuidade técnica mais direta consiste em medir a resposta do motor com acelerômetro fixado à estrutura vestível. Essa instrumentação permitiria relacionar o comando PWM à aceleração, à frequência mecânica e à amplitude do estímulo em condições controladas.

O mapeamento entre \textit{velocity} e PWM também deve ser revisto. Os ensaios mostraram que a referência superior igual a 30 comprime parte relevante da faixa dinâmica do MPD218. A distribuição da faixa completa de 1 a 127, ou uma função não linear ajustada experimentalmente, pode reduzir a saturação e preservar diferenças entre níveis de interação.

Novos ensaios do motor podem empregar ordem randomizada dos valores PWM, maior número de repetições e medição simultânea de tensão, corrente e aceleração. No canal visual, medições fotométricas permitiriam caracterizar a relação entre PWM e brilho e, posteriormente, definir mapeamentos independentes para as modalidades tátil e visual.

A latência fim a fim também pode ser medida com uma referência temporal comum entre o acionamento do pad e as saídas físicas. Esse procedimento permitiria separar os atrasos de aquisição MIDI, processamento em Python, transmissão Wi-Fi, recepção pelo ESP32 e atuação.

Na alimentação, ensaios de autonomia, consumo e temperatura podem orientar a escolha das células e a evolução para uma placa dedicada e uma estrutura mais compacta. A integração mecânica também pode ser refinada para reduzir volume e melhorar o contato do atuador com a pele.

A etapa decisiva para o objetivo de acessibilidade é a avaliação participativa com pessoas surdas ou com deficiência auditiva. Estudos futuros podem investigar discriminação de intensidade, pulsação, ritmo e padrões temporais, além de conforto, preferência e facilidade de aprendizagem. A participação da população-alvo e de profissionais ligados à comunidade surda pode orientar decisões de mapeamento que não podem ser estabelecidas apenas por ensaios técnicos.
